% ==================================================================
%  APPENDIX (unlimited)
% ==================================================================

\appendix

\section{Proof of the Four-Term Identity (Theorem~\ref{thm:four-term})}
\label{app:four-term}

Under Assumption~\ref{ass:packed}, the expected log certified time for
design $(M, q)$ at depth $D$ is:
\begin{equation}
  \E[-\log T] = \sum_s \pi_s \log\bigl(1 - (1 - p(s,M))^D\bigr)
    + \sum_s \pi_s \log q(s).
\end{equation}
The first sum is the competence contribution; the second is the routing
contribution.  The log-performance gap between $(M_0, q_0)$ and $(M, q)$
is:
\begin{align}
  \Delta &= \E[-\log T(M,q)] - \E[-\log T(M_0, q_0)] \nonumber \\
  &= \underbrace{\log\frac{\kappa_0}{\kappa(M)}}_{\text{cost}}
    + \underbrace{\sum_s \pi_s \log\frac{p(s,M)}{p(s,M_0)}}_{\varphi}
    + \underbrace{\sum_s \pi_s \log\frac{q(s)}{q_0(s)}}_{\text{routing difference}}.
\end{align}
The routing difference decomposes via the chain rule for KL divergence.
Define the posterior mode distribution $\pi_D^M(s) \propto \pi_s \cdot
P_D(s,M)$ where $P_D(s,M) = 1 - (1-p(s,M))^D$.  Then:
\begin{align}
  \sum_s \pi_s \log\frac{q(s)}{q_0(s)}
  &= \MI_\pi(S; D_M)
    - \E_D\bigl[\KL(\pi_D^M \| q_D)\bigr]
    + \text{constant w.r.t.\ } q_0 \nonumber \\
  &= G - \mismatch + \text{const}.
\end{align}
Combining yields the four-term identity~\eqref{eq:four-term}. \qed

\section{Lagrangian Non-Connection}
\label{app:lagrangian}

Consider the packed-family optimization with $\Mn{=}2$ modes, two models,
and a cost budget~$B$.  The Lagrangian is:
\[
  \mathcal{L}(q, M, \lambda) = \E[-\log T(M, q)]
    + \lambda\bigl(\kappa(M) \cdot D - B\bigr).
\]
The KKT conditions yield a single scalar multiplier $\lambda^* \ge 0$
as the shadow price of the budget constraint.  The optimal routing satisfies
$q_s^* = \pi_s$ (posterior-matched), and the optimal objective equals
$\Hent(\pi)$, the mode entropy, independent of~$M$.

The four decomposition terms $(\text{cost}, \varphi, G, \mismatch)$ are
components of the primal objective evaluated at any feasible design
point---a zeroth-order algebraic property.  Shadow prices measure
sensitivity to constraint relaxation---a first-order envelope property.
No linear combination of the single shadow price $\lambda^*$ reproduces
the four-term decomposition, except in degenerate cases where the budget
constraint is not binding.

The structural reason is that the decomposition is a \emph{pointwise
identity} valid at every $(M, q)$, not a first-order optimality condition
valid only at the optimum.  The decomposition's value to optimization is
as a separable evaluator, not as a dual decomposition. \qed

\section{Near-Optimality Proof (Proposition~\ref{prop:near-opt})}
\label{app:near-opt}

\begin{proof}
Since $\hat{d}$ minimizes $\hat{f}$ over $\cD$:
\begin{align}
  f(\hat{d})
  &\le \hat{f}(\hat{d}) + \delta
    \tag{$\delta$-accuracy for $\hat{d}$} \\
  &\le \hat{f}(d_{\mathrm{opt}}) + \delta
    \tag{$\hat{d}$ is best estimated} \\
  &\le f(d_{\mathrm{opt}}) + 2\delta.
    \tag{$\delta$-accuracy for $d_{\mathrm{opt}}$}
\end{align}

\textbf{Tightness.}
Take $k{=}2$, $f(d_1) = 0$, $f(d_2) = 2\delta - \epsilon$.
Set $\hat{f}(d_1) = +\delta$, $\hat{f}(d_2) = \delta - \epsilon$.
Algorithm selects $d_2$; regret $= 2\delta - \epsilon \to 2\delta$.

\textbf{Confirmation bound.}
Since $\hat{d} \in \cC$, $\tilde{f}(d^*) \le \tilde{f}(\hat{d})$.
Applying $\eta$-accuracy:
$f(d^*) \le \tilde{f}(d^*) + \eta \le \tilde{f}(\hat{d}) + \eta
\le f(\hat{d}) + 2\eta \le f(d_{\mathrm{opt}}) + 2\delta + 2\eta$. \qed
\end{proof}

\section{Evaluation Complexity Proof}
\label{app:eval-complexity}

Phase~1: $m$ models $\times$ $\Mn$ modes $\times$ $n_0$ instances
$= m \cdot \Mn \cdot n_0$ evaluations.
Phase~2: $k \cdot \Mn$ arithmetic operations, zero evaluations.
Phase~3: $c \cdot N_{\mathrm{conf}}$ evaluations.
Total: $N_{\mathrm{Alg1}} = m \cdot \Mn \cdot n_0 + c \cdot N_{\mathrm{conf}}$.

Racing: $k \cdot n_{\mathrm{eval}}$ evaluations (with
$n_{\mathrm{eval}} = O(\log(k/\beta)/\delta^2)$ for Hoeffding-based
accuracy).

Savings factor:
$N_{\mathrm{race}}/N_{\mathrm{Alg1}} = k \cdot n_{\mathrm{eval}} /
(m \cdot \Mn \cdot n_0 + c \cdot N_{\mathrm{conf}})$.
For fixed $m, \Mn, c, N_{\mathrm{conf}}$, this is $\Theta(k)$.

\textbf{Asymptotic:}
$N_{\mathrm{Alg1}} = O(1)$ in $k$;
$N_{\mathrm{race}} = O(k)$;
$N_{\mathrm{SHA}} = O(B)$ with first-round per-design budget
$O(B/(k\log k)) \to 0$. \qed

\section{Sample Complexity Cascade Proof}
\label{app:cascade}

Chain three results.  From Hoeffding applied to $m \cdot \Mn$ mode--model
pairs with $n_0 = N/(m \cdot \Mn)$ samples each:
\[
  \delta \le \sqrt{\frac{\log(2\Mn/\dconf)}{2n_0}}
  = \sqrt{\frac{m \cdot \Mn \cdot \log(2\Mn/\dconf)}{2N}}.
\]
From Proposition~\ref{prop:near-opt}: regret $\le 2\delta$.
Substituting:
\[
  \text{regret} \le 2\sqrt{\frac{m \cdot \Mn \cdot \log(2\Mn/\dconf)}{2N}}.
\]
Bernstein refinement: replace $1/(2n_0)$ with
$3\pmin/n_0$ when $\pmin \le 1/2$:
\[
  \text{regret}_{\mathrm{Bern}} \le
  2\sqrt{\frac{3m \cdot \Mn \cdot \pmin \cdot \log(2\Mn/\dconf)}{N}}.
\]
Improvement factor: $\sqrt{6\pmin}$; at $\pmin = 0.02$, this is
$\approx 0.35$, a $2.9\times$ tighter bound. \qed

\section{Independence Proofs (Proposition~\ref{prop:independence})}
\label{app:independence}

\subsection{C64: Post-training equivalence does not imply orchestration equivalence}

\textbf{Construction.}
$\Mn{=}2$, $m{=}2$ ($M_{\mathrm{big}}, M_{\mathrm{small}}$).
Shared competence:
$p(1, M_{\mathrm{big}}) = 0.8$, $p(2, M_{\mathrm{big}}) = 0.3$,
$p(1, M_{\mathrm{small}}) = 0.4$, $p(2, M_{\mathrm{small}}) = 0.6$.
Mode distributions: $\pi_A = (0.7, 0.3)$, $\pi_B = (0.3, 0.7)$.

Post-training equivalence holds (identical per-mode, per-model success
probabilities).

Orchestration equivalence fails: $\KL(\pi_A \| \pi_B) = 0.339 > 0$.
The posterior routing distributions at $D{=}1$ are
$\pi_1^A \approx (0.757, 0.243)$ vs.\
$\pi_1^B \approx (0.364, 0.636)$,
with $\KL(\pi_1^A \| \pi_1^B) \approx 0.320$.
The ordering persists at all depths by monotonicity of
$P_D(s) = 1 - (1-p^*(s))^D$.

\subsection{C65: Orchestration equivalence does not imply post-training equivalence}

\textbf{Construction.}
$\pi_C = \pi_D = (0.5, 0.5)$.
$\tau_C$: $p_C(1, M_{\mathrm{big}}) = 0.9$,
  $p_C(2, M_{\mathrm{big}}) = 0.2$,
  $p_C(1, M_{\mathrm{small}}) = 0.3$,
  $p_C(2, M_{\mathrm{small}}) = 0.7$.
$\tau_D$: $p_D(1, M_{\mathrm{big}}) = 0.6$,
  $p_D(2, M_{\mathrm{big}}) = 0.1$,
  $p_D(1, M_{\mathrm{small}}) = 0.2$,
  $p_D(2, M_{\mathrm{small}}) = 0.8$.

Orchestration equivalence holds at the routing-policy level:
$M_{\mathrm{big}}$ wins mode~1 and $M_{\mathrm{small}}$ wins mode~2
in both distributions, at all depths (by strict monotonicity).

Post-training equivalence fails:
$|p_C(1, M_{\mathrm{big}}) - p_D(1, M_{\mathrm{big}})| = 0.3$.
The competence term $\varphi$ differs by $|\varphi_C - \varphi_D|
\approx 0.413$.

\section{Three-Case Transfer Proof (Theorem~\ref{thm:transfer})}
\label{app:transfer}

\textbf{Setup.}
Write the transfer error as:
\[
  \Delta = f(d_S^*; \tau_T) - f(d_T^*; \tau_T)
  \le |f(d_S^*; \tau_T) - f(d_S^*; \tau_S)|
    + |f(d_S^*; \tau_S) - f(d_T^*; \tau_T)|.
\]
Since $d_S^*$ optimizes $\tau_S$: $f(d_S^*; \tau_S) \le f(d_T^*; \tau_S)$.
It suffices to bound $|f(d; \tau_T) - f(d; \tau_S)|$ term by term.

\textbf{Cost term:} $|C_T - C_S| = 0$ (model-dependent, not
distribution-dependent).

\textbf{Competence term:}
$|\varphi_T - \varphi_S| \le 2\alpha_{\mathrm{PT}}$ (each of two
log-terms shifts by at most $\alpha_{\mathrm{PT}}$).

\textbf{Information gain:}
$|G_T - G_S| \le \sqrt{2\alpha_O} \cdot \log\Mn$ by Pinsker's inequality
($\|\pi_S - \pi_T\|_{\mathrm{TV}} \le \sqrt{\alpha_O/2}$) and the
Audenaert entropy-continuity bound.

\textbf{Routing mismatch:}
$|\mismatch_T - \mismatch_S| \le \sqrt{2\alpha_O} \cdot \log\Mn$
by the same Pinsker argument applied to the KL divergence of routing
allocations.

\textbf{Case~1:}
$\Delta \le 2\alpha_{\mathrm{PT}} + \sqrt{2\alpha_O}\cdot\log\Mn$.
Under $\alpha_{\mathrm{PT}} \le \tau_1 = \dtarget/2$ and
$\alpha_O \le \tau_2 = \dtarget^2/(2\log^2\Mn)$:
$\Delta \le \dtarget + \dtarget = 2\dtarget$.

\textbf{Case~2:}
Keep $M_S^*$, re-estimate routing from target pilot.
Residual error: $\alpha_{\mathrm{PT}}$ (competence) $+ 2\dacc$
(re-estimation, by Proposition~\ref{prop:near-opt} applied to the
restricted design space $\{(M_S^*, D) : D \in D_{\mathrm{values}}\}$).
Cost: $\Mn \cdot n_0'$ for one model instead of $m \cdot \Mn \cdot n_0$
for all models---a factor of $m$ savings.

\textbf{Case~3:}
No transfer. Full Algorithm~\ref{alg:decomp-guided} on $\tau_T$.
$\Delta \le 2\dacc$ by Proposition~\ref{prop:near-opt}. \qed

\section{Decision Threshold Derivation}
\label{app:thresholds}

Set each term in the Case~1 bound equal to $\dtarget$:
\begin{align}
  2\alpha_{\mathrm{PT}} &\le \dtarget
    &\implies \tau_1 &= \dtarget/2, \\
  \sqrt{2\alpha_O}\cdot\log\Mn &\le \dtarget
    &\implies \tau_2 &= \dtarget^2/(2\log^2\Mn).
\end{align}
For $\dtarget = 0.1$ and $\Mn = 5$:
$\tau_1 = 0.05$, $\tau_2 \approx 0.0019$.
The orchestration threshold $\tau_2$ is much smaller than $\tau_1$,
reflecting the amplification by $\log\Mn$: even small mode-distribution
shifts can cause significant routing errors.

\section{Transfer Savings Proof (Proposition~\ref{prop:transfer-savings})}
\label{app:transfer-savings}

Case~1: $\mathcal{C}_{\mathrm{target}} = 0$.
Case~2: Pilot evaluates one model on $\Mn$ modes instead of $m$ models.
Pilot cost: $\Mn \cdot n_0$ vs.\ $m \cdot \Mn \cdot n_0$.
Ratio: $1/m$.
Case~3: $\mathcal{C}_{\mathrm{target}} = m \cdot \Mn \cdot n_0 +
c \cdot N_{\mathrm{conf}} = \mathcal{C}_S$. No savings. \qed

\section{Additional Experiment Details}
\label{app:experiments}

\subsection{EXP10 Setup}

Five models (tiny, small, medium, large, xlarge) with costs
$\kappa \in \{0.5, 1, 3, 10, 25\}$.
Per-mode success probabilities scale linearly with model size.
$\Mn{=}5$ modes with Zipf distribution
$\pi \propto (1, 1/2, 1/3, 1/4, 1/5)$.
Cost--quality weight $\mu = 0.3$.
Pilot: $n_0 = 200$ per mode per model.
Confirmation: $c = 5$, $N_{\mathrm{conf}} = 500$.
Racing: $n_{\mathrm{eval}} = 500$ per design.
50 Monte Carlo replications per condition.

\subsection{EXP13 Setup}

Same model and mode parameters as EXP10.
All methods given a matched budget of 7,500 total evaluations.
BO: RBF kernel, length scale 0.5, EI acquisition, 5 initial random
points, 150 evaluations per query.
SHA: budget $= 7{,}500$, $\lceil\log_2 k\rceil$ rounds.
50 MC replications per condition.

\subsection{EXP14 Setup}

Three models (specialist\_A, generalist\_B, specialist\_C) with
heterogeneous competence profiles.
$k = 51$ designs (3 models $\times$ 17 depths).
$\dtarget = 0.15$, yielding $\tau_1 = 0.075$, $\tau_2 \approx 0.0043$.
100 MC replications per grid cell.

\subsection{EXP15 Data Sources}

Semi-synthetic pass rates calibrated to published SWE-bench-lite
results~\citep{jimenez2024swebench}.
Repository-specific pass rates derived from leaderboard entries
(circa early 2025).
Not exact per-instance predictions; structural findings (mode
heterogeneity, absence of crossing structure) are robust to the
exact numbers.
